\section{Introduction}
Economic counterfactuals require more than a numerical solution at a single calibration. A researcher changes an operating benefit, an adjustment cost, or the distribution of a financial shock and asks whether a policy response reflects an economic mechanism rather than approximation error. Repeating an optimization at every parameter value is one response. Reusing an approximate policy is another, but it is useful only when its welfare loss and the scope of its validity can be evaluated. This paper studies that problem for dynamic economies in which control changes both financial resources and the agent's preferences.

Neural Bellman Operators (NBO) separate policy proposal, evaluation, improvement, and certification. A critic represents continuation value, an actor proposes a feasible action, an improvement calculation checks that action, and an independent evaluation measures the induced policy. We use this separation to develop a counterfactual method with an explicit distinction between changes in rewards and changes in transition laws. The certificate, rather than a training objective or a preferred approximation architecture, determines which economic conclusions are warranted.

The first distinction is mathematical. For fixed transitions, the value of a policy is affine in reward coefficients. Policy caches and optimistic linear support exploit precisely this structure. Once a parameter changes the transition law at every date, even a fixed policy can have a nonconvex value in that parameter. We identify the cross-operator term that makes an uncorrected endpoint chord fail and construct a positive backward correction yielding an upper value throughout a parameter interval. Its additional error is second order in interval width under explicit Lipschitz conditions. Independently evaluated policy values have a finite Bernstein representation for an affine mixture of transition laws. Combining their maximum with the corrected upper value and statewise policy switching gives a welfare certificate at every state and date over a continuum of transition counterfactuals. The result also admits approximate neural endpoint values after a certified residual correction.

The second distinction is economic. In a portfolio economy with adjustable risk aversion and first-exit liquidation, a constant operating benefit $m$ and a liquidation increment $\ell$ affect behavior through the contrast $d=m-\rho\ell$. The annuity identity explains which contract changes matter, but not whether a portfolio reversal is caused by preference formation. We partition policies by the sign of their first risky position. The slope of a regular indifference locus with respect to adjustment cost equals the difference in adjustment effort between the two risk classes divided by their difference in discounted duration. A separate option-value decomposition compares deliberate preference adjustment with the otherwise identical economy in which adjustment is unavailable. It identifies when the extra adjustment option, rather than an operating-duration incentive common to both economies, changes the ranking of risky positions.

The finite implementation retains consumption, risky investment, costly preference adjustment, stochastic preference and wealth shocks, their covariance, and the original settlement rule. It includes the union of both historical action meshes and the frozen neural proposals. The reward-transfer bank retains its all-state, all-date $10^{-3}$ welfare certificate and its separate first-action sign exclusions. The new uncertainty experiment changes the probability of two one-step shock laws, each with its own killed transition kernel, and certifies their entire mixture interval. This changes the dynamics; it is not a relabeling of a reward coefficient. The finite-chain mixture is distinguished from changing a Brownian correlation parameter without a discretization bound.

The computational comparisons assign credit to the object doing the work. We execute an exact-oracle scalar optimistic-linear-support comparison, distinguish an initial-distribution objective from simultaneous statewise coverage, and remove all neural actions from the lower policy bank while retaining the original full-menu upper benchmark. In the resource economy, we add an analytically assembled, probability-preconditioned parametric quadratic program that reuses the same structure across fresh queries. It is substantially faster than the learned arm in this short quadratic model. This comparison changes the inference from the earlier crossover: nonlinear continuation can outperform the fitted quadratic approximation without outperforming the exact reusable quadratic program. The transition-law certificate supplies an additional counterfactual capability; it does not depend on suppressing that stronger comparator.

The general error analysis and the broader economic extensions remain part of the argument. A finite-horizon theorem distinguishes evaluation, improvement, boundary, and transition errors. Exact policy improvement provides a conditional benchmark; counterexamples exclude identification of attracting actor stationary points with globally optimal policies. Merton and Epstein--Zin examples test actual evaluation and control updates. Sophisticated temporal selves require a different improvement criterion from ordinary continuation evaluation. Persistent capacity competition tests unilateral dynamic deviations at every date and state. These exercises retain their separate economic purposes rather than becoming an aggregate count of successes.

\subsection{Related Literature and the Contribution}
Policy improvement and approximate dynamic programming precede neural approximation. \citet{jacka2015} study continuous-time improvement under explicit control hypotheses. The nearest neural approaches also combine continuation approximation and policy improvement: fixed-policy PDE evaluation in \citet{kim2025,kim2026}, operator learning in \citet{lee2024}, martingale evaluation in \citet{cai2025}, and direct control approximation in \citet{han2016}. Neural evaluation, separate parameter blocks, and avoidance of a closed-form maximizer are not claims of priority here. Our operator assumptions and error accounts specify when such computations support the economic use of a policy.

Successor features and generalized policy improvement \citep{barreto2017}, policy switching \citep{chang2021}, and policy-cache upper bounds \citep{nemecek2021} provide the reward-transfer foundations. \citet{alegre2022} combine optimistic linear support with successor features and generalized policy improvement. Their work is a necessary comparator: selecting an unsatisfactory reward weight, solving there, and adding its policy is not a new algorithm merely because the reward weight is a contract. We prove below that, with exact statewise endpoint optima and one policy per scalar anchor, our adjacent-policy envelope equals the full scalar-OLS policy envelope. We measure the cost of strengthening an initial-distribution guarantee to simultaneous statewise coverage, but do not present that strengthening alone as a new transfer principle. Our switching rule also differs from full action-value generalized policy improvement; the comparison does not identify the two operators.

The additional transfer result concerns \emph{changing transition laws}, outside the fixed-feature reward correspondence. Lipschitz dependence of value on rewards and transitions is established in the changing-environment literature \citep{csaji2008}. Theorem \ref{thm:corrected_chord} instead gives an explicit endpoint-supersolution correction involving the product of a kernel difference and a continuation-value difference. Its implementation is paired with exact feasible-policy polynomials, a statewise welfare bound, and an independently computed information-relaxation upper value. The latter is an instance of the general approach of \citet{brown2010}, not a new duality theorem. The economic switching-locus result uses the envelope logic of \citet{milgrom2002}; its content is the risk-specific duration--effort comparison and the option-value experiment it determines, not a claim to a new envelope theorem.

The broader computational literature includes recursive and projection methods \citep{stokey1989,judd1998}, sparse grids \citep{brumm2017}, and neural differential-equation solvers \citep{raissi2019,han2018}. Our input-convex critic follows \citet{amos2017}. Riccati recursion and a reusable convex quadratic program are the appropriate structural comparators for the corresponding models. We retain both, including results unfavorable to learned deployment. Thus the paper's incremental case rests on the transition-law upper correction, its operational policy-reuse guarantee, and its use in a mechanism-specific economic analysis, rather than on priority for familiar reward geometry or an unqualified neural speed claim.
